% Part:sets-functions-relations
% Chapter: sets
% Section: schroder-bernstein

\documentclass[../../../include/open-logic-section]{subfiles}

\begin{document}

\olfileid[pt]{sfr}{siz}{sb}

\olsection{A Noção de Tamanho e Schröder-Bernstein}

\begin{explain}
Eis um pensamento intuitivo: se $A$ não é maior do que $B$ e $B$ não é maior do
que $A$, então $A$ e $B$ são equipotentes. Sejamos honestos: se isto estivesse
\emph{errado}, dificilmente poderíamos justificar que a noção definida de
equipotência tem a ver com comparações de ``tamanhos'' entre conjuntos.
Felizmente, o pensamento intuitivo está certo. Isto justifica-se pelo Teorema de
Schr\"oder-Bernstein.
\end{explain}

\begin{thm}[Schr\"oder-Bernstein]
	\ollabel{thm:schroder-bernstein}
	Se $\cardle{A}{B}$ e $\cardle{B}{A}$,
	então $\cardeq{A}{B}$.
\end{thm}

\begin{explain}
Noutras palavras, se existir !!a{injection} de $A$ para~$B$ e !!a{injection} de
$B$ para~$A$, então existe !!a{bijection} de $A$ para~$B$.

Este resultado é, contudo, realmente bastante \emph{difícil} de provar. De
fato, embora Cantor o tivesse enunciado, outros provaram-no.\footnote{Para mais
sobre a história, ver por exemplo \citet[pp.~165--6]{Potter2004}.}
\oliflabeldef{sfr:cardinals:card-sb:sec}{Só estaremos em condições de
\emph{provar} Schr\"oder-Bernstein na \olref[sfr][cardinals][card-sb]{sec}.}{}%
Por agora, pode (e deve) aceitá-lo como verdadeiro.

Felizmente, Schr\"oder-Bernstein é \emph{correto} e legitima pensarmos nas
relações que definimos, O.s.\@, $\cardeq{A}{B}$ e $\cardle{A}{B}$, como tendo a ver
com ``tamanho''. Além disso, Schr\"oder-Bernstein é muito \emph{útil}. Pode ser
difícil imaginar !!a{bijection} entre dois conjuntos equipotentes. O Teorema de
Schr\"oder-Bernstein permite decompor a comparação em casos em que só precisamos
de !!a{injection} do primeiro para o segundo e vice-versa.
\end{explain}
\end{document}
